\documentclass[11pt,letterpaper]{article}
\usepackage[T1]{fontenc}
\usepackage[utf8]{inputenc}
\usepackage{lmodern}
\usepackage[margin=1in,headheight=15pt]{geometry}
\usepackage{microtype}
\usepackage{amsmath,amssymb,amsthm,mathtools}
\usepackage{booktabs,longtable,tabularx,array}
\usepackage{enumitem}
\usepackage{xcolor}
\usepackage{listings}
\usepackage{fancyhdr}
\usepackage{hyperref}
\usepackage{bookmark}
\definecolor{ink}{HTML}{18344B}
\definecolor{accent}{HTML}{176A75}
\definecolor{pale}{HTML}{F3F6F8}
\hypersetup{colorlinks=true,linkcolor=ink,urlcolor=accent,citecolor=accent,
  pdftitle={Lagrange's Four-Square Theorem in Intuitionistic Primitive Recursive Arithmetic},
  pdfauthor={Research report prepared for Vladimir Reshetnikov}}
\pagestyle{fancy}
\fancyhf{}
\fancyhead[L]{\small\color{ink}Four squares in intuitionistic PRA}
\fancyhead[R]{\small 22 November 2026}
\fancyfoot[C]{\thepage}
\setlength{\parskip}{0.42em}
\setlength{\parindent}{0pt}
\setlist{itemsep=0.2em,topsep=0.3em,leftmargin=1.6em}
\setcounter{tocdepth}{2}
\allowdisplaybreaks[2]
\newcommand{\NN}{\mathbb N}
\newcommand{\ZZ}{\mathbb Z}
\newcommand{\ipra}{\mathsf{iPRA}}
\newcommand{\PR}{\mathsf{PR}}
\newcommand{\Nrm}[1]{\mathsf N(#1)}
\newcommand{\monus}{\mathbin{\dot{-}}}
\newcommand{\Prime}{\mathsf{Prime}}
\newcommand{\Init}{\mathsf{Init}}
\newcommand{\Step}{\mathsf{Step}}
\newcommand{\rep}{\mathsf{Rep}}
\newcommand{\abs}[1]{\lvert#1\rvert}
\newcommand{\code}[1]{\texttt{#1}}
\newcommand{\ind}[1]{[\,#1\,]}
\newtheorem{theorem}{Theorem}[section]
\newtheorem{lemma}[theorem]{Lemma}
\newtheorem{proposition}[theorem]{Proposition}
\newtheorem{corollary}[theorem]{Corollary}
\theoremstyle{definition}
\newtheorem{definition}[theorem]{Definition}
\newtheorem{remark}[theorem]{Remark}
\lstset{basicstyle=\ttfamily\small,backgroundcolor=\color{pale},
  frame=single,rulecolor=\color{pale},breaklines=true,columns=fullflexible,
  keepspaces=true,showstringspaces=false,xleftmargin=0.6em,xrightmargin=0.6em,
  aboveskip=0.7em,belowskip=0.7em}
\newcolumntype{Y}{>{\raggedright\arraybackslash}X}
\newcolumntype{L}[1]{>{\raggedright\arraybackslash}p{#1}}

\begin{document}
\hypersetup{pageanchor=false}
\begin{titlepage}
\vspace*{0.55in}
{\small\sffamily\color{accent}CONSTRUCTIVE ARITHMETIC \quad / \quad RESEARCH PACKET}
\vspace{0.25in}

{\Huge\bfseries\color{ink}Lagrange's\par Four-Square Theorem\par}
\vspace{0.18in}
{\Large in intuitionistic primitive recursive arithmetic\par}
\vspace{0.4in}
{\large A complete mathematical argument, explicit witnesses,\par
and a Lean formalization roadmap}
\vspace{0.35in}

\begin{abstract}
Every natural number is the sum of four squares, and this assertion admits
an explicitly witnessed proof in intuitionistic primitive recursive
arithmetic. We construct four primitive recursive functions $W_i$ and prove,
using only quantifier-free induction, that
\[
  W_0(n)^2+W_1(n)^2+W_2(n)^2+W_3(n)^2=n.
\]
The construction consists of a bounded residue search for each prime,
an explicit descent with separate even and odd multiplier branches, and a
bounded trial-factor loop for arbitrary inputs. The central descent is
iterated a prescribed number of times; all its invariants are arithmetic
equalities and decidable inequalities. Integer and tuple notation is reduced
to primitive recursive coding over the natural numbers.

The packet includes an exact-arithmetic Python implementation, reproducible
checks, uncompiled Lean source outlining the complete construction, and
pinned ProveIt interface references. The mathematical proof is supplied in
full. The Lean arithmetic obligations and the final syntactic iPRA derivation
remain to be mechanized; no completed machine verification is claimed.
\end{abstract}
\vfill
\begin{tabular}{@{}ll@{}}
Prepared for & Vladimir Reshetnikov\\
Date & 22 November 2026\\
Foundational target & Intuitionistic first-order PRA, quantifier-free induction\\
Proof route & Bounded modular seed, parity-first descent, bounded factor loop
\end{tabular}
\end{titlepage}
\hypersetup{pageanchor=true}
\pagenumbering{roman}
\begingroup
\small
\setlength{\parskip}{0pt}
\tableofcontents
\endgroup
\clearpage
\pagenumbering{arabic}

\section{The theorem and the precise foundational target}
\label{sec:foundations}

\begin{theorem}[Witnessed four-square theorem in $\ipra$]
\label{thm:main}
There are four explicitly specified primitive recursive functions
$W_i:\NN\to\NN$, $0\le i<4$, such that $\ipra$ proves
\begin{equation}
\label{eq:main}
 \forall n\;\bigl(W_0(n)^2+W_1(n)^2+W_2(n)^2+W_3(n)^2=n\bigr).
\end{equation}
It also proves $W_i(n)\le n$ for each $i$. Consequently it proves
\begin{equation}
\label{eq:exists}
 \forall n\,\exists a,b,c,d\le n\quad n=a^2+b^2+c^2+d^2.
\end{equation}
The convention is $\NN=\{0,1,2,\ldots\}$.
\end{theorem}

The functions will be defined before their correctness is concluded. Thus
the assertion is stronger than a pointwise existence argument: the witnesses
are supplied by fixed arithmetic programs, and the proof establishes their
correctness uniformly inside the specified theory.

\subsection{What is meant by intuitionistic PRA}

We use the standard first-order presentation of $\ipra$. Its only sort is
$\NN$. Its language has zero, successor, and a function symbol for every
primitive recursive definition, generated from projections by composition
and primitive recursion. If
\[
 f(0,\vec x)=g(\vec x),\qquad
 f(t+1,\vec x)=h(t,f(t,\vec x),\vec x),
\]
the corresponding defining equations are axioms. The non-logical axioms
also state that successor is injective and never zero, and include induction
\[
 \bigl(A(0,\vec u)\land
       \forall t(A(t,\vec u)\to A(t+1,\vec u))\bigr)
 \to \forall t\,A(t,\vec u)
\]
only when $A$ is quantifier-free. Free parameters are allowed.

The logic is intuitionistic first-order logic with equality: the usual
introduction and elimination rules for $\land,\lor,\to,\forall,\exists$,
falsehood elimination, and equality reflexivity and substitution. No general
law of excluded middle, choice principle, or double-negation elimination is
added. This is the presentation implemented by ProveIt's
\code{ConciseIPRA.PR} and \code{PRAAx}~\cite{proveit}.

Writing a new primitive recursive function in this presentation adds a
name for a definition, not a stronger existence axiom. A smaller fixed
language may instead compile the definitions into its established primitive
recursive term language.

\subsection{Decidable arithmetic and bounded statements}

Truncated subtraction $a\monus b=\max(a-b,0)$, equality tests, order tests,
parity, natural quotient and remainder, bounded sums, and bounded products
are primitive recursive. Their elementary laws have quantifier-free
induction proofs. The details of quotient/remainder and coding used here
are given in Appendix~\ref{app:coding}.

Every bounded predicate occurring in this report has a specified primitive
recursive characteristic bit. For a bit-valued primitive recursive
predicate $b(i,\vec u)$, define
\[
 \operatorname{All}(0,\vec u)=1,\qquad
 \operatorname{All}(k+1,\vec u)
   =\operatorname{All}(k,\vec u)b(k,\vec u).
\]
The equation $\operatorname{All}(t,\vec u)=1$ expresses the finite conjunction
$\bigwedge_{i<t}b(i,\vec u)=1$. A similar fold expresses finite disjunction.
Negation of a bit is $1\monus b$; conjunction is multiplication; disjunction
can be written $\min(1,b+c)$. Thus the finite tests are decidable in $\ipra$.

When a proof below inducts on a statement displayed with bounded
quantifiers, the actual induction formula is its characteristic-bit
equation. This is a concrete elimination of bounded quantifiers, not an
extension of the induction scheme to arbitrary quantified formulas.
Appendix~\ref{app:folds} gives the needed fold and search reasoning.

In particular, eliminating a failed bounded search is constructive. A
computed search flag is either $0$ or $1$. If the $0$ branch entails a
contradiction, the $1$ branch remains and contains actual recorded witnesses.
This is a decidable case split, not an application of double-negation
elimination to an arbitrary existential assertion.

\subsection{Signed integers and finite records are notation}

For readability we calculate with $\ZZ$ and four-tuples. Internally, a signed
integer is a natural code, and a tuple is a nested primitive recursive pair.
All operations used below, including signed Euclidean remainder and exact
division by a positive natural, have primitive recursive implementations.
Appendix~\ref{app:coding} gives an explicit encoding and operations.

For $x=(x_0,x_1,x_2,x_3)\in\ZZ^4$, put
\[
 \Nrm{x}=x_0^2+x_1^2+x_2^2+x_3^2.
\]
This is a natural-valued primitive recursive function of the tuple code.
The state assertion $\Nrm{x}=mp$ is therefore quantifier-free after coding.
No quantification over arbitrary functions, sets of integers, or infinite
sequences is needed anywhere in the proof.

\section{The algebraic identity that carries the construction}
\label{sec:identity}

\begin{definition}[Euler product]
For $x=(a,b,c,d)$ and $y=(e,f,g,h)$, define
\begin{equation}
\label{eq:product}
 x*y=\bigl(
 ae-bf-cg-dh,\;
 af+be+ch-dg,\;
 ag-bh+ce+df,\;
 ah+bg-cf+de
 \bigr).
\end{equation}
Write $\bar y=(e,-f,-g,-h)$.
\end{definition}

Only this explicit polynomial operation will be used. No theory of
quaternion algebras is required.

\begin{lemma}[Euler's four-square identity]
\label{lem:euler}
For all integer four-tuples $x,y$,
\[
 \Nrm{x*y}=\Nrm{x}\Nrm{y},\qquad
 x*\bar x=(\Nrm{x},0,0,0),\qquad
 \Nrm{\bar x}=\Nrm{x}.
\]
These are finite polynomial identities and have $\ipra$ proofs.
\end{lemma}

\begin{proof}
The four coordinates of $x*y$ are the rows of the following matrix applied
to $(e,f,g,h)^{\mathsf T}$:
\[
 L_x=\begin{pmatrix}
 a&-b&-c&-d\\
 b&a&-d&c\\
 c&d&a&-b\\
 d&-c&b&a
 \end{pmatrix}.
\]
Every column has squared length $a^2+b^2+c^2+d^2$. The six distinct
column dot products are
\begin{align*}
 -ab+ab+cd-dc&=0,&
 -ac-bd+ac+bd&=0,&
 -ad+bc-bc+ad&=0,\\
 bc-ad+ad-bc&=0,&
 bd+ac-bd-ac&=0,&
 cd-dc-ab+ab&=0.
\end{align*}
Expansion of the squared norm therefore has coefficient $\Nrm{x}$ on each
of $e^2,f^2,g^2,h^2$ and zero on every mixed term. This proves the first
identity. Substitution of $(a,-b,-c,-d)$ into~\eqref{eq:product} gives the
second, and changing signs leaves squares unchanged, giving the third.

The displayed matrix is only a compact arrangement of sixteen integer
expressions. Each calculation reduces to associativity, commutativity, and
distributivity in coded integer arithmetic; it invokes no linear-algebraic
existence theorem. Those ring laws themselves follow from elementary
quantifier-free arithmetic induction. In ProveIt the signed-pair identity
can be produced by \code{ZT.derive} and \code{ipra\_ring}; see
Section~\ref{sec:formalization}.
\end{proof}

The key consequence is constructive multiplicative closure: given actual
tuples of norms $r$ and $s$, formula~\eqref{eq:product} supplies an actual
tuple of norm $rs$. Taking absolute values gives nonnegative coordinates
without changing the norm.

\section{Bounded elementary number theory}
\label{sec:elementary}

\begin{definition}[Primality]
For a natural number $p$, $\Prime(p)$ means
\[
 p\ge2\quad\text{and}\quad
 \forall d<p\;\bigl(2\le d\to p\bmod d\ne0\bigr).
\]
This notation abbreviates a primitive recursive Boolean test. Divisibility
by a positive natural is also a remainder-zero test and supplies the
quotient as its witness.
\end{definition}

\begin{lemma}[A bounded least divisor is prime]
\label{lem:lpf}
For $q\ge2$, scan the finite list $2,3,\ldots,q$ and let $p$ be its first
divisor of $q$. This search succeeds, $2\le p\le q$, $p\mid q$, and $p$ is
prime. All the functions and correctness assertions are available in $\ipra$.
\end{lemma}

\begin{proof}
The last candidate $q$ divides $q$, so the scan cannot fail. If
$2\le d<p$ and $d\mid p$, then $d\mid q$: multiply the exact quotient
identities for $p/d$ and $q/p$. That contradicts the recorded failure of
every candidate smaller than $p$. Each candidate is a decidable test, and
the scan's correctness is proved by quantifier-free induction on its index.
No assertion about existence or uniqueness of a full prime factorization
is used.
\end{proof}

\begin{lemma}[Prime divisibility of a product]
\label{lem:euclid}
If $p$ is prime and $p\mid uv$ for signed integers $u,v$, then
$p\mid u$ or $p\mid v$. This disjunction has a constructive $\ipra$ proof.
\end{lemma}

\begin{proof}
Compute $a=u\bmod p$, so $0\le a<p$. If $a=0$, the first disjunct holds.
Assume $a>0$. We give the bounded arithmetic proof that $a$ has a
B\'ezout relation with $p$.

Initialize remainders and signed coefficients by
\[
 r_0=p=1p+0a,\qquad r_1=a=0p+1a.
\]
At an active stage, with $r_1>0$, compute
$Q=r_0\mathbin{\mathrm{div}}r_1$ and $r_2=r_0-Qr_1$ and replace
$(r_0,r_1)$ by $(r_1,r_2)$. Replace their coefficient pairs in exactly the
same way. If $r_1=0$, retain the entire record. Run exactly $p+1$ stages.

At every active stage $0\le r_2<r_1$; hence the second remainder decreases
by at least one until it is zero. Quantifier-free induction gives the bound
$r_1(t)\le a\monus t$. After $p+1$ stages it is zero. The first remainder
stays positive, and both remainders stay at most $p$. The coefficient
identities are preserved because the same linear operation is applied to
both sides of each identity.

For every $1\le d\le p$, the relation
\[
 (d\mid r_0\land d\mid r_1)
 \quad\Longleftrightarrow\quad
 (d\mid r_1\land d\mid r_0-Qr_1)
\]
follows by substitution of the quotient identities, in each direction.
Consequently the preserved common-divisor invariant is the finite conjunction
\[
 \bigwedge_{1\le d\le p}
 \bigl((d\mid r_0(t)\land d\mid r_1(t))
       \leftrightarrow(d\mid p\land d\mid a)\bigr).
\]
Represent this conjunction by its Boolean fold. It holds initially, and
the displayed one-step equivalence proves its preservation by transitivity,
using quantifier-free induction. At the final record $(g,0)$,
it follows, by taking $d=g$, that $g\mid p$ and $g\mid a$.
Here $1\le g\le p$.

Primality gives $g=1$ or $g=p$: a value $2\le g<p$ is excluded directly
by the primality test. The latter alternative is impossible because
$0<a<p$. Thus $g=1$. The computed coefficient record supplies actual
integers $s,t$ with
\[
 sp+ta=1.
\]
The coefficient update is primitive recursive. For an optional explicit
size bound, every quotient is at most $p$, and the maximum absolute
coefficient grows by a factor of at most $p+1$ per stage. Thus
$(p+1)^{p+1}$ bounds its size after $p+1$ stages.

Because $u\equiv a\pmod p$, the hypothesis gives $p\mid av$.
Multiplying $sp+ta=1$ by $v$ shows that $p\mid v$. The initial test
$a=0$ decides which disjunct is obtained. Every loop has a prescribed
bound and a coded quantifier-free invariant, so the proof uses only
allowed induction.
\end{proof}

\section{A finite modular seed for every prime}
\label{sec:seed}

The prime $2$ already has the tuple $(1,1,0,0)$. A prime greater than $2$
is odd, since an even number greater than $2$ has the forbidden divisor
$2$. Hence write an odd prime as $p=2h+1$, with $h\ge1$.

\begin{lemma}[Two short lists of square residues]
\label{lem:residue}
For $p=2h+1$ prime, there are explicitly computable $a,b\le h$ with
\[
 a^2+b^2+1\equiv0\pmod p.
\]
One may find them by at most $(h+1)^2$ tests.
\end{lemma}

\begin{proof}
For $0\le i,j\le h$, form the residues in $\{0,\ldots,p-1\}$
\[
 A_i=i^2\bmod p,\qquad B_j=(-1-j^2)\bmod p.
\]
Each list is injective. If $A_i=A_j$, then
$p\mid(i-j)(i+j)$. Lemma~\ref{lem:euclid} gives $p\mid i-j$ or
$p\mid i+j$. The bounds $\abs{i-j}<p$ and $0\le i+j\le2h=p-1$
then give $i=j$ in the first case, or $i=j=0$ in the second.
Indeed an integer divisible by $p$ and of absolute value less than $p$
must be zero, by the exact quotient equation. Equality of two $B$ entries
reduces to the same square-residue argument.

Scan the $(h+1)^2$ pairs in lexicographic order, storing the first pair
with $A_a=B_b$ and a hit flag. We prove that the failure flag is impossible
by an explicit finite counting argument. If there were no cross-list
match, the concatenation $L$ of the two lists would have
$2(h+1)=p+1$ pairwise distinct entries, each smaller than $p$.

For $r<p$, let
\[
 C(r)=\sum_{k=0}^{p}\ind{L_k=r},
\]
where the bracket is the equality-test bit. Because the entries are
pairwise distinct, at most one summand is $1$, so $C(r)\le1$.
For each $k$, exactly one of the $p$ equality tests $\ind{L_k=r}$ is $1$,
because $L_k<p$. Counting the finite incidences in the two orders gives
\begin{equation}
\label{eq:count}
 p+1
 =\sum_{k=0}^{p}\sum_{r=0}^{p-1}\ind{L_k=r}
 =\sum_{r=0}^{p-1}C(r)
 \le p,
\end{equation}
a contradiction. The interchange of the two finite sums is proved by
successively adding rows and distributing addition. These are nested
quantifier-free inductions on the two finite bounds. The assertion that
at most one summand is $1$ can either be recorded as a Boolean prefix
condition or proved by induction on the count; no set cardinalities or
infinite combinatorial principles are involved.

The hit flag is a computed bit. Its $0$ case has just been excluded, so it
is $1$ and its stored pair satisfies $A_a=B_b$. This equality is precisely
the desired congruence. The scan's index $k<(h+1)^2$ can be decoded as
\[
 a=k\mathbin{\mathrm{div}}(h+1),\qquad b=k\bmod(h+1),
\]
which specifies the search as a single primitive recursive bounded loop.
\end{proof}

\begin{proposition}[Small positive multiple of a prime]
\label{prop:seed}
For each prime $p$ a primitive recursive seed construction returns
$m_0\in\NN$ and $x^{(0)}\in\ZZ^4$ with
\[
 1\le m_0<p,\qquad \Nrm{x^{(0)}}=m_0p.
\]
For odd $p$ it can ensure $m_0\le(p-1)/2$.
\end{proposition}

\begin{proof}
For $p=2$ take $m_0=1$ and $x^{(0)}=(1,1,0,0)$.
For odd $p=2h+1$, use the successful pair of Lemma~\ref{lem:residue} and set
\[
 x^{(0)}=(a,b,1,0),\qquad
 m_0=(a^2+b^2+1)\mathbin{\mathrm{div}}p.
\]
The division is exact. Its numerator is positive, so $m_0\ge1$.
Since $h\ge1$,
\[
 a^2+b^2+1\le2h^2+1\le2h^2+h=hp.
\]
Cancellation of the positive factor $p$ gives $m_0\le h<p$.
For nonprime inputs the same finite program is made total by returning
an arbitrary default, namely $(1,(0,0,0,0))$, when no pair is found.
The theorem asserts correctness only under the decidable prime condition.
\end{proof}

\section{Explicit descent of the multiplier}
\label{sec:descent}

Fix a prime $p$. A valid state is a pair $s=(m,x)$ satisfying
\begin{equation}
\label{eq:valid}
 I_p(s):\quad 1\le m<p\quad\land\quad \Nrm{x}=mp.
\end{equation}
The task is to reach $m=1$ while preserving this assertion. Instead of
choosing a least possible multiplier, we define the next state directly.

\subsection{Even multipliers: exact parity halving}

\begin{lemma}[Halving an even norm]
\label{lem:half}
There is a primitive recursive function $H:\ZZ^4\to\ZZ^4$ such that
\[
 2\mid\Nrm{x}\quad\Longrightarrow\quad 2\Nrm{H(x)}=\Nrm{x}.
\]
\end{lemma}

\begin{proof}
Since $z^2\equiv z\pmod2$, an even norm means that the number of odd
coordinates is even. For $x=(a,b,c,d)$ choose one of three orders:
\[
 (a,b,c,d)\quad\hbox{if }a\equiv b\pmod2;
\]
otherwise choose $(a,c,b,d)$ if $a\equiv c\pmod2$, and $(a,d,b,c)$ in the
remaining case. These are explicit parity tests.

In the first two cases, one chosen pair has equal parity, and the parity of
the total norm forces the remaining pair to have equal parity as well.
In the last case $b$ and $c$ both have parity opposite to $a$, so they
have equal parity. The parity of the total then forces $d$ to have the
same parity as $a$. Hence each chosen pair has equal parity in all cases.

Write the reordered tuple as $(u,v,w,z)$ and define
\begin{equation}
\label{eq:half}
 H(x)=\left(\frac{u+v}{2},\frac{u-v}{2},
              \frac{w+z}{2},\frac{w-z}{2}\right).
\end{equation}
All four divisions are exact, including those of negative numerators.
For either pair,
\[
 (u+v)^2+(u-v)^2=2u^2+2v^2.
\]
Substitute the exact quotient identities, add the two pair identities,
and cancel the positive integer $2$. The result is
$2\Nrm{H(x)}=u^2+v^2+w^2+z^2=\Nrm{x}$.
\end{proof}

If a valid state has even $m>1$, set
\[
 m'=m/2,\qquad x'=H(x).
\]
Then $1\le m'<m$, and Lemma~\ref{lem:half} gives $\Nrm{x'}=m'p$.
Thus this branch strictly decreases the multiplier.

\subsection{Odd multipliers: centered residues and exact division}

Assume now that $m>1$ is odd. For an integer $z$, let $t=z\bmod m$,
where $0\le t<m$, and define
\begin{equation}
\label{eq:center}
 \rho_m(z)=
 \begin{cases}
 t,&2t\le m,\\
 t-m,&2t>m.
 \end{cases}
\end{equation}
Because $m$ is odd, this gives
\begin{equation}
\label{eq:centerbound}
 \rho_m(z)\equiv z\pmod m,\qquad
 \abs{\rho_m(z)}\le\frac{m-1}{2}.
\end{equation}
Put $r_i=\rho_m(x_i)$ and $r=(r_0,r_1,r_2,r_3)$.

\begin{lemma}[The new multiplier]
\label{lem:newmult}
For a valid state with odd $m>1$, the quotient
\[
 k=\Nrm{r}\mathbin{\mathrm{div}}m
\]
is exact and satisfies $1\le k<m$.
\end{lemma}

\begin{proof}
Congruence of the coordinates implies congruence of their squares and
their sums. Therefore
\[
 \Nrm{r}\equiv\Nrm{x}=mp\equiv0\pmod m,
\]
so $\Nrm{r}=mk$. Moreover,
\[
 mk=\Nrm{r}
 \le4\left(\frac{m-1}{2}\right)^2
 =(m-1)^2<m^2.
\]
Cancellation of $m>0$ gives $k<m$. The displayed estimate can be written
without fractions by setting $m=2h+1$: it is simply
$\Nrm{r}\le4h^2<(2h+1)^2$.

It remains to rule out zero. If $k=0$, the nonnegative squares in
$\Nrm{r}=0$ must all be zero. Thus all $r_i=0$, and $m$ divides every
$x_i$. Compute the exact quotients $w_i=x_i/m$. Then
\[
 m^2\Nrm{w}=\Nrm{x}=mp,
\]
and cancellation gives $m\Nrm{w}=p$. This would make $m$ a divisor of
$p$ with $1<m<p$, contradicting the primality test. Since the natural
number $k$ is decidably either zero or positive, it follows that $k\ge1$.
\end{proof}

\begin{lemma}[The new representing tuple]
\label{lem:quotient}
Under the assumptions of Lemma~\ref{lem:newmult}, every coordinate of
$z=x*\bar r$ is divisible by $m$, and the exact quotient tuple
$x'=z/m$ satisfies $\Nrm{x'}=kp$.
\end{lemma}

\begin{proof}
For clarity, with $x=(a,b,c,d)$ and $r=(r,s,t,u)$, the four numerators are
\begin{align*}
 z_0&= ar+bs+ct+du,\\
 z_1&=-as+br-cu+dt,\\
 z_2&=-at+bu+cr-ds,\\
 z_3&=-au-bt+cs+dr.
\end{align*}
Modulo $m$ one can replace $a,b,c,d$ by $r,s,t,u$, respectively.
The first expression becomes $r^2+s^2+t^2+u^2=\Nrm{r}=mk$;
each of the other three becomes zero by cancellation of paired products.
Thus $m\mid z_i$ for each $i$, with exact signed quotient $x'_i$.

Euler's identity and invariance of the norm under conjugation give
\[
 m^2\Nrm{x'}=\Nrm{z}
 =\Nrm{x}\Nrm{\bar r}
 =(mp)(mk)=m^2pk.
\]
Cancel the positive integer $m^2$ to obtain $\Nrm{x'}=pk$.
\end{proof}

\subsection{The total transition and its decrease theorem}

The transition on coded states is now completely specified:
\begin{equation}
\label{eq:step}
 \Step(m,x)=
 \begin{cases}
 (m,x),&m\le1,\\
 (m/2,H(x)),&m>1\text{ even},\\
 (\Nrm{r}/m,(x*\bar r)/m),&m>1\text{ odd},
 \end{cases}
\end{equation}
where $r$ is given by~\eqref{eq:center}. Quotient operations are defined
on all inputs; their exactness is a proved property of valid states.
Therefore this is a total primitive recursive function even on malformed
state codes. It does not need a partial ``divide if divisible'' operation.

\begin{proposition}[One-step descent]
\label{prop:step}
If $p$ is prime and $I_p(m,x)$ holds, then $I_p(\Step(m,x))$ holds.
If $m>1$, the new multiplier $m'$ satisfies $1\le m'<m$.
If $m=1$, the state is unchanged.
\end{proposition}

\begin{proof}
The terminal case follows from the first branch of~\eqref{eq:step}.
The even case follows from Lemma~\ref{lem:half}, with $m'=m/2$.
The odd case follows from Lemmas~\ref{lem:newmult} and~\ref{lem:quotient}.
These exhaustive cases are determined by natural-number comparison and
parity, both decidable arithmetic tests.
\end{proof}

\begin{remark}[Why the parity branch matters]
A centered-residue bound for an arbitrary even modulus gives only
$\abs{r_i}\le m/2$, hence $\Nrm{r}/m\le m$. That estimate alone does
not prove strict descent. Treating even multipliers by exact halving removes
the boundary case altogether. For the remaining odd multipliers,
$\abs{r_i}\le(m-1)/2$ yields strict decrease directly. The separate
argument excluding $k=0$ is also essential.
\end{remark}

\section{Termination by a fixed primitive recursive clock}
\label{sec:clock}

We record the precise induction used instead of an appeal to infinite
descent.

\begin{lemma}[An absorbing bounded descent]
\label{lem:clock}
Let $D$ be a primitive recursive operation on natural state codes, $J$ a
quantifier-free state predicate, and $R$ a natural-valued primitive
recursive rank. Suppose arithmetic proves
\[
 J(s)\to J(D(s)),\qquad
 J(s)\to R(D(s))\le R(s)\monus1.
\]
Define $D^0(s)=s$ and $D^{t+1}(s)=D(D^t(s))$ by primitive recursion.
Then quantifier-free induction proves
\begin{equation}
\label{eq:rankinvariant}
 J(s)\to
 \bigl(J(D^t(s))\land R(D^t(s))\le R(s)\monus t\bigr).
\end{equation}
In particular, $R(D^t(s))=0$ whenever $t\ge R(s)$ and $J(s)$ holds.
\end{lemma}

\begin{proof}
For $t=0$, both claims follow from $D^0(s)=s$ and $R(s)\monus0=R(s)$.
Assume the displayed invariant at $t$. Preservation gives $J(D^{t+1}(s))$,
and the rank inequality gives
\[
 \begin{split}
 R(D^{t+1}(s))
 &\le R(D^t(s))\monus1\\
 &\le(R(s)\monus t)\monus1
   =R(s)\monus(t+1).
 \end{split}
\]
Monotonicity and the last identity are elementary laws of truncated
subtraction, proved by quantifier-free induction. Formula
\eqref{eq:rankinvariant}, with state code $s$ as a parameter, is itself
quantifier-free. At $t\ge R(s)$ its right-hand bound is zero.
\end{proof}

For the prime descent, use
\[
 J(s)=I_p(s),\qquad R(m,x)=m\monus1.
\]
If $m=1$, the rank is zero and remains zero. If $m>1$, strict multiplier
decrease gives the required saturated decrease of the rank. Define
\begin{equation}
\label{eq:primeprog}
 s_t(p)=\Step^{\,t}(m_0(p),x^{(0)}(p)),\qquad
 P(p)=\text{the tuple component of }s_p(p).
\end{equation}
Here the clock is the given input $p$, not the first time a test succeeds.

\begin{theorem}[Prime witness]
\label{thm:prime}
The function $P$ is primitive recursive, and $\ipra$ proves
\[
 \Prime(p)\to\Nrm{P(p)}=p.
\]
\end{theorem}

\begin{proof}
The seed is valid by Proposition~\ref{prop:seed}. Apply
Lemma~\ref{lem:clock} and Proposition~\ref{prop:step} to obtain
\[
 I_p(s_t(p)),\qquad
 (m_t\monus1)\le(m_0\monus1)\monus t.
\]
These assertions are obtained by quantifier-free induction, with the
prime-test equation as a parameter assumption. Since $m_0<p$, at $t=p$
the rank is zero. Validity also gives $m_p\ge1$, so $m_p=1$. The norm
part of validity is now $\Nrm{P(p)}=p$.

The seed, transition, and iteration have already been defined as total
primitive recursive functions. The proof does not establish their
termination by a separate unbounded search; it establishes correctness
after a predetermined number of steps.
\end{proof}

\subsection{Two fully worked descent traces}

For $p=7$, the lexicographic seed search returns $(a,b)=(2,3)$, so
\[
 2\cdot7=2^2+3^2+1^2+0^2.
\]
The tuple is reordered as $(2,0,3,1)$ into equal-parity pairs. Halving
sums and differences gives $(1,1,2,1)$ and
$7=1^2+1^2+2^2+1^2$.

For $p=23$, the seed is $(2,8,1,0)$ with multiplier $3$. Centering modulo
$3$ gives $r=(-1,-1,1,0)$ and $\Nrm{r}=3$. Hence $k=1$.
The conjugate product is
\[
 (2,8,1,0)*(-1,1,-1,0)=(-9,-6,-3,-9).
\]
Divide by $3$ to obtain $(-3,-2,-1,-3)$, of norm $23$.

A trace exercising both branches is
\[
 \begin{array}{c|c|c}
 p=31&m&x\\\hline
 \text{seed}&6&(4,13,1,0)\\
 \text{even halving}&3&(2,2,7,6)\\
 \text{odd descent}&1&(1,2,-1,-5).
 \end{array}
\]
Every row satisfies $\Nrm{x}=31m$.

\section{From prime witnesses to every natural number}
\label{sec:assembly}

We now avoid induction on the quantified statement ``$n$ has a
four-square representation.'' A second explicit bounded loop carries
the tuple that is being constructed.

For $n>0$, start with
\[
 q_0=n,\qquad A_0=(1,0,0,0).
\]
For a state $(q,A)$ define
\begin{equation}
\label{eq:factorstep}
 T(q,A)=
 \begin{cases}
 (q,A),&q\le1,\\
 (q/p,\;A*P(p)),&q>1,
 \end{cases}
\end{equation}
where $p$ in the second branch is the first divisor of $q$ in
$2,\ldots,q$. This is the bounded search of Lemma~\ref{lem:lpf}.
Iterate $T$ exactly $n$ times.

\begin{proposition}[Factor-loop invariant]
\label{prop:factor}
For each $n>0$, quantifier-free induction proves
\begin{equation}
\label{eq:factorinvariant}
 1\le q_t,\qquad q_t\Nrm{A_t}=n,\qquad
 (q_t\monus1)\le(n\monus1)\monus t.
\end{equation}
Consequently $q_n=1$ and $\Nrm{A_n}=n$.
\end{proposition}

\begin{proof}
The initial norm is $1$, so the invariant holds at $t=0$.
If $q=1$, the state is unchanged. If $q>1$, Lemma~\ref{lem:lpf}
supplies a prime $p$ dividing $q$, with $2\le p\le q$.
The exact quotient $q'=q/p$ is positive and satisfies $q'<q$:
from $q=pq'$ and $p\ge2$ one obtains $q\ge2q'>q'$.

Theorem~\ref{thm:prime} and Euler's identity give
\[
 q'\Nrm{A*P(p)}
 =q'\Nrm{A}\Nrm{P(p)}
 =q'\Nrm{A}p
 =q\Nrm{A}=n.
\]
Thus the first two parts of the invariant are preserved, and the rank
$q\monus1$ decreases in the saturated sense of Lemma~\ref{lem:clock}.
That lemma supplies the third part. At $t=n$ its bound is zero,
so positivity gives $q_n=1$, after which the norm equation gives
$\Nrm{A_n}=n$.
\end{proof}

\begin{proof}[Proof of Theorem~\ref{thm:main}]
For $n=0$ define $W_i(0)=0$. For $n>0$ take the tuple $A_n$ produced
by~\eqref{eq:factorstep} and define
\[
 W_i(n)=\abs{(A_n)_i}\qquad(0\le i<4).
\]
All components are primitive recursive, by composition of the bounded
searches, coded tuple operations, and the two fixed-clock recursions.
Proposition~\ref{prop:factor} proves that their squares sum to $n$.
This also holds at zero by definition. The split $n=0$ or $n>0$ is
a decidable arithmetic split, giving~\eqref{eq:main} in $\ipra$.

For each coordinate, $W_i(n)^2\le n$ because the other summands are
nonnegative. If $W_i(n)=0$ its bound is immediate. Otherwise write
$W_i(n)=u+1$; then $(u+1)^2\ge u+1$, so $W_i(n)\le n$.
Finally introduce the four existential quantifiers using the four
specific terms $W_i(n)$. This derives~\eqref{eq:exists} intuitionistically.
\end{proof}

\subsection{A compact program specification}

The mathematical program is the following. Its arguments, records, signed
operations, and parity tests are all primitive recursive after the coding
of Appendix~\ref{app:coding}.

\begin{lstlisting}
prime_rep(p):
    if p = 2: return (1,1,0,0)
    if p < 3: return (0,0,0,0)
    h := (p-1) div 2
    search a,b in 0..h for (a*a+b*b+1) mod p = 0
    if the bounded search fails: return (0,0,0,0)
    x := (a,b,1,0); m := norm(x) div p
    repeat exactly p times:
        if m <= 1: retain (m,x)
        else if m is even:
            x := parity_halve(x); m := m div 2
        else:
            r := centered_residues(x,m)
            z := x * conjugate(r)
            (m,x) := (norm(r) div m, z div m)
    return x

four_squares(n):
    if n = 0: return (0,0,0,0)
    q := n; A := (1,0,0,0)
    repeat exactly n times:
        if q = 1: retain (q,A)
        else:
            p := first divisor of q in 2..q
            (q,A) := (q div p, A * prime_rep(p))
    return componentwise_abs(A)
\end{lstlisting}

Simultaneous assignment in this specification uses the old state on every
right-hand side. The distributed Python implementation skips repeated
terminal states by default; its \code{literal\_clock=True} option executes
all the prescribed iterations. Both variants implement the same function
on natural inputs, because the skipped states are absorbing.

\subsection{A still simpler canonical witness definition}

The constructive proof also justifies a very small, though slower, witness
program. Search the $(n+1)^4$ codes for four digits in base $n+1$ and
return the first $(a,b,c,d)$ with $a^2+b^2+c^2+d^2=n$.
Theorem~\ref{thm:main} supplies a successful candidate, since all four
coordinates are at most $n$. Hence the failure sentinel is impossible.
This gives four primitive recursive witnesses defined using natural
arithmetic alone. The descent construction is what proves the bounded
search always succeeds; merely writing the bounded search would not by
itself prove the theorem.

\section{Why the whole argument belongs to \texorpdfstring{$\ipra$}{iPRA}}
\label{sec:audit}

The arithmetic proof above is already constructive. The remaining
foundational question is the strength of its induction. The following
table records every type of induction that carries substantive work.

\begin{longtable}{@{}L{0.22\textwidth}L{0.29\textwidth}L{0.43\textwidth}@{}}
\toprule
Task & Recursion bound & Quantifier-free invariant or equation\\
\midrule
\endhead
Quotient and remainder & Dividend $a$ & $a=mQ+R$ and $R<m$ for $m>0$.\\
Finite search & Prescribed candidate bound & Stored success is a genuine hit;
all earlier failures are a Boolean fold.\\
B\'ezout computation & $p+1$ & Two coefficient identities, positive first
remainder, decreasing second remainder, finite common-divisor bit.\\
Residue counting & $p+1$ rows, $p$ columns & Equations of finite sums and
prefix counts; range and distinctness use finite Boolean folds.\\
Prime descent & $p$ & $I_p(s_t)$ and
$(m_t\monus1)\le(m_0\monus1)\monus t$.\\
General assembly & Original $n$ & $q_t\ge1$, $q_t\Nrm{A_t}=n$, and
$(q_t\monus1)\le(n\monus1)\monus t$.\\
\bottomrule
\end{longtable}

Each state field is a primitive recursive term in the loop counter and
the input parameters. Consequently the norm equalities and rank
inequalities in the last two rows are literally quantifier-free after
expanding the state projections. The prime condition is a fixed
primitive recursive bit equation. Its presence as a parameter assumption
does not add a quantifier to the induction formula.

The finite counting argument may be displayed in ordinary mathematical
notation using statements about all entries of a list. Its actual
implementation uses bounded sums of equality-test bits. Their laws are
proved by induction on numerical prefix lengths. This is a finite
combinatorial proof over coded data, requiring no comprehension axiom.

All uses of contradiction concern a decidable branch already under
consideration: the failed search flag, a zero new multiplier, or a
forbidden divisor. General classical reductio is unnecessary. All
witnesses used after such a split are either already stored in a record
or obtained by exact arithmetic. There is no application of countable
choice, Markov's principle, unbounded minimization, or induction on an
existential representation predicate.

\subsection{Explicit bounds on intermediate arithmetic}

The bound $m<p$ gives $\Nrm{x}=mp<p^2$, hence $\abs{x_i}<p$ throughout
prime descent. In its odd branch, $\abs{r_i}<p$ and each raw coordinate
of $x*\bar r$ has absolute value less than $4p^2$. In its even branch,
raw sums and differences have absolute value less than $2p$.
After division, the norm invariant again bounds the new coordinates by
$p$. The seed search has at most $((p+1)/2)^2$ candidates for odd $p$.

In the outer loop, $q\Nrm{A}=n$ and $q\ge1$ give
$\Nrm{A}\le n$, so every accumulated coordinate has absolute value at
most $n$. Every selected prime is at most $q\le n$.
Before normalization, a product coordinate is bounded in absolute value
by $4n^2$. A fixed nesting of polynomial-size pair codes remains
primitive recursive. These elementary register bounds are useful for
mechanization, though ordinary primitive recursion does not require
polynomial bounds on its intermediate values.

The construction is intended to expose logical strength, not to optimize
running time in the binary length of $n$. Under a unit-cost convention for
integer arithmetic and division, the direct fixed-clock program has a
crude $O((n+1)^3)$ bound on its arithmetic operations: at most $n$ outer
stages, each with a search of at most $(n+1)^2$ candidates and at most $n$
descent steps. This observation is not a polynomial-time claim in
$\log n$ and is not needed for the proof.

\section{Lean formalization: definitions, obligations, and object proofs}
\label{sec:formalization}

\subsection{The exact status of the supplied Lean files}

The packet contains two Lean source outlines:
\begin{itemize}
\item \code{lean/FourSquares.lean}: total executable definitions and proposed
proof scripts assembling the theorem from six explicit arithmetic contracts.
It imports Lean's standard library and the \code{omega} tactic, without
importing Mathlib's four-square theorem.
\item \code{lean/ProveItTarget.lean}: the exact first-order iPRA goal in
ProveIt's syntax, a stronger goal parameterized by four unary primitive
recursive symbols, and a proposed object-level Euler-identity proof fragment.
This file requires the full pinned ProveIt workspace.
\end{itemize}

\textbf{Neither file was elaborated or compiled in this environment: a Lean
compiler was not available.} No completed Lean four-square proof or final
$\mathsf{PRAAx}$ derivation is claimed. The files expose the work still to
be done as named local obligations; they contain no admitted theorem
declarations. Absence of an admission is not a compilation receipt, and
a conditional theorem is not an unconditional proof of its conclusion.

The preceding sections prove the mathematical content of the obligations.
The separate formalization ledger, \code{OBLIGATIONS.md}, decomposes the
remaining work into 31 concrete steps, including the final syntax-level
translation.

\subsection{Definitions and their correspondence to the proof}

The host Lean development uses an explicit record with four integer
coordinates, and two records for the loops. The complete definitions are in
the file, with this correspondence:

\begin{longtable}{@{}L{0.42\textwidth}L{0.52\textwidth}@{}}
\toprule
Mathematical object & Lean identifier\\
\midrule
\endhead
Finite conjunction and bounded first hit & \code{allLt}, \code{firstLt}\\
Primality bit and first divisor & \code{primeB}, \code{IsPrime}, \code{leastDivisor}\\
Integer four-tuple, norm, product, conjugate &
\code{Z4}, \code{Z4.normSq}, \code{Z4.mul}, \code{Z4.conj}\\
Centered residues & \code{centered}, \code{centered4}\\
Three pairings and exact halving & \code{halfPairs}, \code{evenHalf}\\
Initial multiple of a prime & \code{seedState}\\
Transition~\eqref{eq:step} & \code{descentStep}\\
State assertion~\eqref{eq:valid}, rank & \code{PrimeState}, \code{descentRank}\\
Primitive recursive iteration & \code{iterate}\\
Prime witness~\eqref{eq:primeprog} & \code{primeRep}\\
Outer factor transition & \code{factorStep}\\
Outer norm invariant and rank & \code{FactorInvariant}, \code{factorRank}\\
Signed result and natural result & \code{signedFourSquares}, \code{fourSquares}\\
Unfinished arithmetic interface & \code{ArithmeticContracts}\\
Conditional correctness and existence & \code{fourSquares\_correct},
\code{lagrange\_from\_contracts}\\
\bottomrule
\end{longtable}

\par\noindent
\begin{minipage}{\linewidth}
For example, the recursion is structurally on its explicit clock:
\begin{lstlisting}
def iterate (step : A -> A) : Nat -> A -> A
  | 0, s     => s
  | k + 1, s => step (iterate step k s)
\end{lstlisting}
\end{minipage}\par
Here \code{A} is a host-level type parameter. The actual iPRA instances
use natural state codes and fixed first-order primitive recursive operations.
The type-polymorphic presentation does not imply that iPRA itself has
higher-order objects.

\subsection{The six local contracts}

The file's \code{ArithmeticContracts} structure contains exactly the
following mathematical claims:
\begin{enumerate}
\item \textbf{Seed correctness:} a prime input yields a valid seed state.
This is Proposition~\ref{prop:seed}, including finite-search correctness.
\item \textbf{Descent preservation:} $I_p(s)$ implies $I_p(\Step(s))$.
This is Proposition~\ref{prop:step}.
\item \textbf{Descent rank decrease:}
$R(\Step(s))\le R(s)\monus1$ for $R=m\monus1$ on valid states.
The terminal, even, and odd branches were proved separately.
\item \textbf{Factor-loop preservation:} assuming the proved prime witness
property, the outer transition preserves
$q\ge1$ and $q\Nrm{A}=n$.
\item \textbf{Factor-loop rank decrease:} on this invariant,
$(q'\monus1)\le(q\monus1)\monus1$.
\item \textbf{Naturalization:} replacing each integer by its natural
absolute value preserves the sum of squares.
\end{enumerate}

The source proposes generic preservation and rank-bound proofs, then uses
these six fields to assemble
\begin{lstlisting}
theorem fourSquares_correct (C : ArithmeticContracts) (n : Nat) :
    (fourSquares n).normSq = n

theorem lagrange_from_contracts (C : ArithmeticContracts) (n : Nat) :
    exists a b c d : Nat, n = a*a+b*b+c*c+d*d
\end{lstlisting}
The actual Lean files use standard Lean existential notation. The argument
\code{C} is an explicit hypothesis record, and no instance of it is asserted.
To finish the host formalization, prove the listed elementary lemmas,
construct that record, and specialize these statements to it.

Useful local proof methods are elementary rewriting for quotient equations,
\code{omega} for linear inequalities and truncated subtraction, and a
polynomial normalization tactic for the two Euler identities and the
halving identity. Nonlinear inequalities such as
$(m-1)^2<m^2$ should be reduced to explicit polynomial expansions plus
positive-factor reasoning; they should not be silently given to an
unrestricted high-level theorem.

\subsection{The required distinction: host Lean truth versus iPRA derivability}

A theorem of type
\[
 \forall n:\NN,\;\exists a,b,c,d:\NN,\quad n=a^2+b^2+c^2+d^2
\]
is a host Lean theorem about natural numbers. On its own, it does not
state that iPRA proves the sentence. A complete machine formalization of
the foundational claim must construct an object derivation of type
\[
 ([]:\operatorname{List}(\mathsf{Fm}\ \mathsf{PTm}\ 0))
 \ \vdash[\mathsf{PRAAx}]\ \mathsf{fourSquareSentence}.
\]
This type is explicitly defined as \code{FourSquareObjectGoal} in the
second source file; no inhabitant is postulated.

The innermost matrix has five variables. Inside
$\forall n\,\exists a\,\exists b\,\exists c\,\exists d$, ProveIt's
de Bruijn indices are
\[
 0=d,\quad1=c,\quad2=b,\quad3=a,\quad4=n.
\]
Thus the matrix is
\[
 v_4\doteq(v_3v_3+v_2v_2)+v_1v_1+v_0v_0,
\]
with the association fixed in \code{natSquareSum}. The four existential
constructors are wrapped inside one universal constructor. This specifies
the target without any ambiguity about which witness each binder denotes.

For the stronger result, the source defines a family parameter
$W:\operatorname{Fin}4\to\mathsf{PR}\,1$ and the closed object sentence
\[
 \forall n\;\left(n=\sum_{i<4}W_i(n)W_i(n)\right).
\]
The remaining construction is to compile the four actual programs into
four \code{PR 1} terms and produce the derivation of that identity.
Existential introductions then derive the ordinary sentence.

Host-level generic induction lemmas may quantify over arbitrary predicates.
Their use here is legitimate for the iPRA audit only because the actual
instantiations are the coded quantifier-free predicates in
Section~\ref{sec:audit}. A machine proof must record these QF certificates
and invoke \code{PRA.qfInd} or an already proved restricted induction
interface. Merely observing that Lean accepted a host induction proof
would not establish this restriction.

\subsection{Concrete integration with the inspected ProveIt snapshot}

The source review is pinned to commit
\begin{center}
\texttt{3a521d0265941f9e9d4ad450818365a78a6c9a73}.
\end{center}
At that snapshot the root toolchain specifies Lean \code{v4.32.0};
the dependency documentation records Mathlib commit
\texttt{81a5d257c8e410db227a6665ed08f64fea08e997}~\cite{proveit}.
The following actual source interfaces supply relevant object derivations:

\begin{longtable}{@{}L{0.34\textwidth}L{0.60\textwidth}@{}}
\toprule
Source module & Integration use\\
\midrule
\endhead
\code{Core.Definitions} & \code{PTm.toPR}, \code{recDef},
\code{rec\_zero}, \code{rec\_succ} compile the fixed arithmetic programs
and expose their recursion equations.\\
\code{Core.Search} & \code{Search.search\_spec} and
\code{Search.search\_found} give bounded-search success, minimality,
and failure information. Its search returns the bound on failure.\\
\code{Core.Pigeonhole} & \code{Pigeonhole.pigeonhole} supplies explicit
colliding indices for a coded list of $K+1$ values below $K$. Its import
closure is substantial; the report also proves the simple counting argument
directly.\\
\code{Core.EuclidLemma} & \code{euclid\_lemma\_at} supplies the prime
divisor disjunction as an object derivation.\\
\code{AnalyticNumberTheory.Primes} & \code{lpf\_dvd} and
\code{lpf\_prime} justify a bounded least-prime-factor search.\\
\code{Core.Pairing} & \code{pairT}, \code{fst\_pair},
\code{snd\_pair}, and \code{simF} support state coding and simultaneous
primitive recursion.\\
\code{Core.Integers} & \code{ZT} represents differences of natural terms;
\code{ZT.derive} and \code{ipra\_ring} can produce the fixed integer
polynomial identities as object proofs.\\
\code{Core.ClockedDescent} & \code{clocked\_at} and \code{clocked}
give the QF refutational descent principle. The witness-producing absorbing
loop here should instead use its displayed rank invariant directly.\\
\bottomrule
\end{longtable}

These are source-inspected interfaces, not newly recompiled dependencies.
The copied reference files retain their source identifiers and the
repository's MIT-0 license. They are a reference subset, not a transitive
import closure. Build the object target in the full pinned repository.

There is an important encoding distinction. ProveIt's \code{ZT} is a pair
$(u,v)$ representing $u-v$, with no canonicality condition. Equality of
$(u,v)$ and $(r,s)$ is the natural equation $u+s=r+v$. The canonical
integer code in Appendix~\ref{app:coding} has an explicit PR bridge to
this representation: normalize a difference pair to
$(u\monus v,v\monus u)$, then encode its sign and nonzero magnitude.
The polynomial identity can be proved in \code{ZT} first and transported
through this bridge. Thus no assumption about quotient-type integer
objects is needed inside iPRA.

\code{Core.ClockedDescent} has a complete explicit proof in the inspected
source, and its introducing commit reports a successful build and axiom
audit. This packet did not independently recompile it, and no independent
CI check result was retrieved~\cite{clocked}. Moreover, that interface
forbids a nonempty QF state predicate closed under strict descent; it does
not already supply the tuple-producing absorbing iteration of this report.
Keeping that difference explicit prevents a refutational theorem from
being mistaken for a finished witness constructor.

A sensible new module family is
\code{IPRADerivations/DiophantineNumberTheory/FourSquares/}, with modules
\code{Identity}, \code{Seed}, \code{Halving}, \code{OddDescent},
\code{PrimeWitness}, \code{Assembly}, and \code{Statement}.
Their dependency order is the order of this report. The six local
contracts and the QF invariants specify the precise mathematical boundary
between them.

\subsection{Relation to the existing ordinary Lean theorem}

Mathlib already provides \code{Nat.sum\_four\_squares}, together with the
Euler identity and a prime case. The inspected source uses
\code{classical}, \code{Nat.findX}, and induction on multiplicative
factorization~\cite{mathlib}. This is useful as a comparison implementation.
It is not used as a premise here: the explicit bounded witnesses and QF
induction audit are what establish the requested weak-theory result.

\begin{samepage}
The inspected ProveIt file is
\begin{quote}\small
\path{Analysis/FabiusFunction/Lean/FabiusFunction/FourSquareTheorem.lean}
\end{quote}
It uses Mathlib's theorem to obtain an ordinary existence corollary. Its
source is distinct from a $\mathsf{PRAAx}$ derivation~\cite{proveit}.
The source inventory found no four-square-named module in the inspected
iPRA derivation tree; this is a filename observation, not a claim that
no related material occurs anywhere in the repository.
\end{samepage}

\section{Executable artifacts and verification evidence}
\label{sec:validation}

\subsection{Files and reproducibility}

\begin{longtable}{@{}L{0.38\textwidth}L{0.56\textwidth}@{}}
\toprule
Artifact & Purpose\\
\midrule
\endhead
\code{report.tex}, \code{sections/*.tex}, \code{report.pdf} & Complete
mathematical proof and foundational audit, with the source for rebuilding
this PDF.\\
\code{code/four\_squares.py} & Exact-integer implementation of the bounded
seed, parity halving, odd descent, and factor accumulation.\\
\code{code/validate.py} & Reproducible symbolic and finite checks.\\
\code{validation/results.json} & Machine-readable results and exact
test scope.\\
\code{validation/prime\_traces.json} & Explicit descent traces, including
negative coordinates and both transition branches.\\
\code{lean/FourSquares.lean} & Uncompiled algorithm definitions and
conditional host-theorem assembly.\\
\code{lean/ProveItTarget.lean} & Uncompiled exact object sentences and
proposed Euler-identity fragment for the pinned repository.\\
\code{lean/OBLIGATIONS.md} & Named remaining formalization steps.\\
\code{references/proveit/} & Pinned reference sources, interface excerpts,
license, and provenance manifest.\\
\code{validation/proof\_audit.txt}, \code{validation/code\_review.txt} &
Independent mathematical and implementation review notes, with stated scope.\\
\bottomrule
\end{longtable}

The constructor itself uses only the Python standard library. The symbolic
validation script additionally uses SymPy. From the packet directory:
\begin{lstlisting}
python3 code/four_squares.py 2026
python3 code/four_squares.py 31 --prime-trace
python3 code/four_squares.py 128 --literal-clock
python3 code/validate.py
make pdf
\end{lstlisting}
The recorded validation used Python 3.12.14 and SymPy 1.14.0.

\subsection{Checks that were actually run}

The following checks passed and are recorded in
\code{validation/results.json}:
\begin{itemize}
\item Exact symbolic expansion of Euler's identity returned zero remainder.
The conjugate-product diagonal identity and the halving polynomial identity
were also checked symbolically.
\item Every signed ordered four-tuple satisfying
$\Nrm{x}=mp$ with $p\in\{2,3,5,7,11,13\}$ and $1\le m<p$ was tested.
This gives \textbf{19,992 states}: 376 terminal states, 10,992 even-halving
states, and 8,624 odd-centered states. All passed norm preservation,
positive new multiplier, strict active decrease, and, where applicable,
exact coordinate divisibility.
\item The complete constructor was checked for every
$0\le n\le1000$ (1,001 inputs). Each returned four nonnegative integers
whose squares sum exactly to $n$.
\item All 168 primes at most 1,000 were checked for seed correctness and
prime-representation correctness. The active traces exercised 129 even
steps and 194 odd steps, with at most five active steps in that range.
Primality classifications were compared with an independent finite sieve.
\item The early-exit implementation and the literal fixed-clock version
agreed for $0\le n\le128$ and for primes at most 257. An independent
bounded two-pair enumerator also produced valid representations through 128.
\item Additional exact outputs were checked for 2,026, 1,000,000, and
20,261,122, as recorded below.
\end{itemize}

\begin{center}
\begin{tabular}{r@{\qquad}rrrr}
\toprule
$n$&$W_0(n)$&$W_1(n)$&$W_2(n)$&$W_3(n)$\\
\midrule
0&0&0&0&0\\
1&1&0&0&0\\
7&1&1&2&1\\
23&3&2&1&3\\
31&1&2&1&5\\
65&3&2&4&6\\
310&0&6&15&7\\
2,026&0&0&1&45\\
1,000,000&0&1,000&0&0\\
20,261,122&3,706&926&1,157&2,081\\
\bottomrule
\end{tabular}
\end{center}

The finite tests check implementation behavior, especially sign conventions,
parity pairings, exact division, and clock handling. The universal theorem
is established by the proof and QF induction audit, not inferred from the
test range. Symbolic expansion verifies a fixed polynomial identity but
does not produce a Lean or iPRA proof object.

\subsection{Completion boundary}

The human-readable constructive argument and explicit primitive recursive
witness construction are complete. The Python implementation and the
reported checks are complete. The PDF is compiled from the distributed
LaTeX source.

The host Lean local arithmetic contracts remain to be proved and compiled.
The natural-code correspondence and the final closed
$\mathsf{PRAAx}$ derivation remain to be assembled in ProveIt. The supplied
source outlines, exact object goals, pinned interfaces, and obligation
ledger describe those tasks concretely. They do not turn the remaining
formalization into an already verified artifact.

\appendix
\section{Explicit primitive recursive coding and search invariants}
\label{app:coding}

This appendix supplies the numerical details behind the record notation.
It also explains why the displayed finite logical arguments do not hide
stronger induction.

\subsection{Natural quotient and remainder}

Fix $m>0$. Define the pair $(Q_m(a),R_m(a))$ by recursion on $a$:
\[
 (Q_m(0),R_m(0))=(0,0).
\]
If the current remainder is $R$, set
\[
 (Q',R')=
 \begin{cases}
 (Q+1,0),&R+1=m,\\
 (Q,R+1),&R+1<m.
 \end{cases}
\]
For a total definition, the second branch can be the default branch;
the invariant ensures its displayed inequality. Simultaneous recursion is
implemented by pairing. Quantifier-free induction proves
\[
 a=mQ_m(a)+R_m(a),\qquad 0\le R_m(a)<m.
\]
The update proves the identity directly in each case. Uniqueness follows
because two distinct quotients differ by at least one, while two
remainders in $[0,m)$ differ in absolute value by less than $m$.
Set $Q_0(a)=0$ and $R_0(a)=a$ as total defaults. This establishes the
primitive recursive division and remainder operations used in the proof.

Positive-factor cancellation is elementary arithmetic: if $c>0$ and
$ca=cb$, then $a<b$ would imply $cb\ge ca+c>ca$, and $b<a$ would
give the symmetric contradiction. Decidability of natural order leaves
$a=b$. The multiplication-order laws used here follow by quantifier-free
induction from addition and multiplication definitions.

\subsection{A canonical integer code}

Encode a signed integer by
\[
 \operatorname{enc}(z)=
 \begin{cases}
 2z,&z\ge0,\\
 2(-z)-1,&z<0.
 \end{cases}
\]
This display describes the intended meaning. The internal functions are
defined directly on natural codes. For a code $c$, put
\[
 P(c)=\begin{cases}c/2,&c\text{ even},\\0,&c\text{ odd},\end{cases}
 \qquad
 M(c)=\begin{cases}0,&c\text{ even},\\(c+1)/2,&c\text{ odd}.\end{cases}
\]
Thus $c$ represents $P(c)-M(c)$ and its magnitude is
$\operatorname{mag}(c)=(c+1)\mathbin{\mathrm{div}}2$.
At most one of $P(c),M(c)$ is nonzero.

For natural $u,v$, define the code of their difference by
\[
 E(u,v)=
 \begin{cases}
 2(u\monus v),&v\le u,\\
 2(v\monus u)\monus1,&u<v.
 \end{cases}
\]
Then canonical coded addition, negation, and multiplication are
\begin{align*}
 c\oplus d&=E(P(c)+P(d),M(c)+M(d)),\\
 \operatorname{neg}(c)&=E(M(c),P(c)),\\
 c\otimes d&=E(P(c)P(d)+M(c)M(d),\;
                 P(c)M(d)+M(c)P(d)).
\end{align*}
These definitions use only natural arithmetic and decidable cases. Their
interpretation as the desired integer operations follows by expansion of
differences. Sign cases reduce their laws to natural polynomial equations.
The square magnitude $\operatorname{mag}(c)^2$ is the natural norm
contribution of the coded integer.

For division by $m>0$, first compute natural magnitude quotient and
remainder $a=mq+r$, $0\le r<m$. A nonnegative code gives quotient $q$
and remainder $r$. For the code of $-a$, use
\[
 \begin{array}{c|cc}
 &\text{signed quotient}&\text{nonnegative remainder}\\\hline
 r=0&-q&0\\
 r>0&-(q+1)&m-r.
 \end{array}
\]
Both rows satisfy $-a=mQ+R$ and $0\le R<m$. In the exact-division case
this is the usual signed exact quotient. For zero divisor choose quotient
zero and remainder equal to the input signed integer. The valid descent
branches never divide by zero.

All signs in the human proof can therefore be translated into ordinary
natural terms. The reference implementation's signed integers and Lean's
\code{Int} are convenient host representations of these functions, not
an additional sort in the object theory.

\subsection{Tuple coding and simultaneous recursion}

A convenient pairing, matching the inspected ProveIt interface, is
\[
 \pi(u,v)=(u+v)^2+u.
\]
It has primitive recursive projections on its image. Given a code $c$,
let $s$ be the largest $j\le c$ with $j^2\le c$, computed by a bounded
scan; put $u=c\monus s^2$ and $v=s\monus u$.
For $c=\pi(u,v)$, write $s_0=u+v$. Since
\[
 s_0^2\le c=s_0^2+u\le s_0^2+s_0<(s_0+1)^2,
\]
the computed square-root floor is $s_0$, and both projections recover
the original $u,v$. On codes outside the image their values are still
defined, and no inverse law is needed there.

Nested pairing codes a signed four-tuple, a multiplier with a four-tuple,
or a remainder with an accumulated tuple. A single primitive recursion
on the resulting code implements every simultaneous update. The defining
equations for the projections then recover the displayed record equations.
There is no recursor on an infinite stream and no need for a sequence
comprehension principle.

\subsection{Bounded search as a primitive recursive function}
\label{app:folds}

Let $b(i,\vec u)\in\{0,1\}$ be a \emph{fixed} primitive recursive test.
The following recursion returns the first successful index below $B$,
or $B$ if there is none:
\begin{align*}
 s(0,\vec u)&=0,\\
 s(k+1,\vec u)&=
 \begin{cases}
 s(k,\vec u),&s(k,\vec u)<k,\\
 k,&s(k,\vec u)=k\text{ and }b(k,\vec u)=1,\\
 k+1,&\text{otherwise}.
 \end{cases}
\end{align*}
The defining cases are decided by natural tests. Quantifier-free induction
proves the following conjunction, using Boolean folds for its bounded parts:
\begin{enumerate}
\item $s(k,\vec u)\le k$;
\item if $s(k,\vec u)<k$, then $b(s(k,\vec u),\vec u)=1$;
\item every $i<s(k,\vec u)$ has $b(i,\vec u)=0$;
\item if $s(k,\vec u)=k$, every $i<k$ has $b(i,\vec u)=0$.
\end{enumerate}
At a successor stage an existing hit is preserved. If there is none,
the new candidate either becomes the stored hit or extends the failed
prefix by one. These are exactly the three branches above. Thus a known
successful candidate $j<B$ forces $s(B,\vec u)<B$, and a non-failure
output gives a genuine witness.

For the finite conjunction fold, the elementary facts needed are prefix
extension and bounded extraction. Extension is the defining equation
\[
 \operatorname{All}(k+1)=\operatorname{All}(k)b(k).
\]
If the product is $1$, both bit factors are $1$. Conversely, if both are
$1$, the product is $1$. For extraction at a fixed parameter $j$, induct
on the quantifier-free formula
\[
 j<k\ \land\ \operatorname{All}(k)=1\quad\to\quad b(j)=1.
\]
In the successor step, decide whether $j<k$ or $j=k$ and use the old
prefix or the last factor. This supplies the usual elimination rule for
the displayed bounded universal statement without quantified induction.
To build a finite conjunction from already available pointwise proofs,
induct on its bit equation, restricting the index to the chosen total
bound by a quantifier-free guard.

The same ideas justify the finite incidence counts in~\eqref{eq:count}.
For an explicitly coded injective list, a prefix contains at most one
entry equal to a fixed residue. The prefix-count invariant is a natural
inequality; absence of repeated equalities is a bounded Boolean condition.
Adding the next equality-test bit preserves the inequality, because a
previous hit and a new hit would be a repeated entry. Finite double-sum
interchange follows by induction on the row count, and within the step
by induction on the column count using associativity of addition.

Here $b$ is a metamathematical schema for each particular test actually
used in the finite program. The object theory has the corresponding
fixed function symbols. Nothing in the proof assumes a primitive
recursive universal evaluator for arbitrary primitive recursive program
codes. Such an assumption is unnecessary and would be an invalid shortcut.

\subsection{A dependency account for the final derivation}

The proof can be assembled in the following finite sequence:
\begin{enumerate}
\item Derive natural arithmetic, coded integer operations, pairing,
finite sums, and bounded-search specifications.
\item Derive the Euler polynomial identity, bounded least-divisor facts,
and the bounded B\'ezout lemma.
\item Prove injectivity of the two residue lists, the finite incidence
contradiction, and the seed's norm and size bounds.
\item Prove the parity-halving identity and the centered odd-descent lemmas,
including positivity and exact quotient divisibility.
\item Apply QF induction to the explicit prime-state and rank invariant
at clock $p$; obtain the prime witness identity.
\item Apply QF induction to the factor-state and rank invariant at clock
$n$; take magnitudes and obtain the four unary witness identities.
\item Introduce the four existential witnesses and universally close the
free input variable using intuitionistic logical rules.
\end{enumerate}

Every mathematical dependency in this list has been proved in the text
or reduced here to its explicit elementary quantifier-free induction.
The corresponding proof objects are the remaining Lean engineering task.


\begin{thebibliography}{9}
\raggedright
\bibitem{proveit}
Vladimir Reshetnikov and ProveIt contributors.
\emph{ProveIt: Logic/PrimitiveRecursiveArithmetic}.
Repository snapshot
\texttt{3a521d0265941f9e9d4ad450818365a78a6c9a73}, inspected 22 November 2026.
\url{https://github.com/VladimirReshetnikov/ProveIt/tree/3a521d0265941f9e9d4ad450818365a78a6c9a73/Logic/PrimitiveRecursiveArithmetic}.
The packet's \code{references/proveit/SOURCE\_MANIFEST.json} records precise
paths, Git blob identifiers, and SHA-256 hashes.

\bibitem{clocked}
ProveIt contributors. \emph{Core/ClockedDescent.lean}.
Git blob \texttt{4b41c93d5788694b9989a83c49d1c8acdaf918f6} at the snapshot above.
See the pinned URL in the source manifest. Its introducing commit is
\texttt{e1d9051b368ebe351cbbc50932982b06f852c03c}.

\bibitem{mathlib}
Chris Hughes and Mathlib contributors.
\emph{Mathlib.NumberTheory.SumFourSquares}.
Source inspected 22 November 2026, Git blob
\texttt{863f3f96bb65314912188d69c38535c394612813}.
\url{https://github.com/leanprover-community/mathlib4/blob/master/Mathlib/NumberTheory/SumFourSquares.lean}.
Documentation:
\url{https://leanprover-community.github.io/mathlib4_docs/Mathlib/NumberTheory/SumFourSquares.html}.

\bibitem{requested}
\emph{Lagrange's four-square theorem}. Wikipedia, user-supplied topic reference,
consulted 22 November 2026.
\url{https://en.wikipedia.org/wiki/Lagrange%27s_four-square_theorem}.
The proof in this report is developed explicitly; this page is not used as
a formal foundational dependency.
\end{thebibliography}
\end{document}
